\documentclass[conference,twocolumn,spanish]{IEEEtran}

\usepackage[utf8]{inputenc}
\usepackage[T1]{fontenc}
\usepackage[spanish,es-tabla]{babel}
\usepackage{booktabs}
\usepackage{tabularx}
\usepackage{pgfplots}
\pgfplotsset{compat=1.18}
\usepackage{hyperref}

\newcommand{\graficaavanceproyecto}[3]{%
\begin{figure}[htbp]
\centering
\begin{tikzpicture}
\begin{axis}[
    xbar,
    width=0.92\columnwidth,
    height=3.6cm,
    bar width=7pt,
    xlabel={Porcentaje de avance (\%)},
    xmin=0,
    xmax=114,
    xtick={0,25,50,75,100},
    symbolic y coords={#2},
    ytick=data,
    tick label style={font=\scriptsize},
    label style={font=\scriptsize},
    xmajorgrids=true,
    grid style={dashed, gray!35},
    nodes near coords={\pgfmathprintnumber{\pgfplotspointmeta}\%},
    every node near coord/.append style={font=\tiny, anchor=west}
]
\addplot[fill=green!55!black, draw=green!70!black] coordinates {#3};
\end{axis}
\end{tikzpicture}
\caption{Porcentaje de avance por hitos: #1.}
\end{figure}
}

\title{Nombre del Proyecto CIAI}

\author{
  \IEEEauthorblockN{Nombre Apellido 1, Nombre Apellido 2}
  \IEEEauthorblockA{
    Centro de Inteligencia Artificial e Innovación (CIAI) -- UTS\\
    Enlace del proyecto: \url{https://github.com/usuario/proyecto-o-pagina}
  }
}

\begin{document}

\maketitle

\section{Descripción del proyecto}
Redacte aquí en un párrafo breve (80--120 palabras) el objetivo del proyecto, el problema que resuelve, las tecnologías o modelos de IA empleados y su impacto esperado en la UTS o la región.

\section{Avances y rúbrica de seguimiento}
Resuma el avance logrado durante la semana y complete la tabla con los criterios definidos por el equipo:

\begin{table}[htbp]
\centering
\caption{Rúbrica de avance del proyecto.}
\label{tab:rubrica_nombre_proyecto}
\scriptsize
\begin{tabularx}{\columnwidth}{>{\raggedright\arraybackslash}p{2.2cm} >{\raggedright\arraybackslash}X c}
\toprule
\textbf{Criterio / Hito} & \textbf{Evidencia observable de cumplimiento} & \textbf{Estado} \\
\midrule
Datos & Dataset limpio, etiquetado y verificado al 100\,\%. & Completado \\
Modelo & Prototipo entrenado y evaluado sobre conjunto de prueba. & En progreso \\
Interfaz & API o aplicación web funcional para demostración. & Pendiente \\
Validación & Pruebas con usuarios piloto y documentación final. & Pendiente \\
\bottomrule
\end{tabularx}
\end{table}

\graficaavanceproyecto{Nombre Corto del Proyecto}{Validación,Interfaz,Modelo,Datos}{
    (Validación, 0)
    (Interfaz, 20)
    (Modelo, 75)
    (Datos, 100)
}

\end{document}
